\subsection{Блок-схема численного метода Эйлера}

На вход реализуемого метода подаются следующие данные:

\begin{itemize}
  \item \(f\) --- номер дифференциального уравнения; целое число.
  \item \(\varepsilon\) --- требуемая точность вычислений, то есть максимально допустимая разность между двумя итерациями; вещественное число двойной точности.
  \item \(a, b\) --- границы интервала вычислений; вещественные числа двойной точности.
  \item \(y_a\) --- начальное условие в точке \(a\); вещественное число двойной точности.
  \item \(M\) --- максимально допустимое количество разбиений сетки; целое число.
\end{itemize}

На выходе метода ожидается вещественное число двойной точности --- приближённое значение \(y(b)\). 

Остановка итерационного процесса происходит при достижении заданной точности \(\varepsilon\) или при достижении лимита разбиений \(M\). Шаг сетки на каждой итерации уменьшается вдвое за счёт удвоения количества узлов \(n\).

Далее приводятся блок-схемы, которые реализуют вспомогательный алгоритм расчета по сетке и основной численный метод.

\vspace*{\fill}

\begin{center}
  \begin{tikzpicture}[gost flowchart]
    \begin{scope}[start chain=flow going below, node distance=8mm]
      \node [data, on chain, join] (input) {Input: $f, y_a, a, b, n$};
      \node [terminator, left= of input] (start) {Start};
      
      \node [process, on chain, join] (calc_h) {$h \leftarrow (b - a) / n$};
      \node [preparation, on chain, join] (init) {
        $\text{func} \leftarrow \text{getFunc}(f)$ \\
        $y \leftarrow y_a$ \\ 
        $x \leftarrow a$ \\ 
        $i \leftarrow 1$
      };
      
      \node [decision, on chain, join] (check_i) {$i \le n$?};
      
      \node [process, on chain] (loop_body) {
        $y \leftarrow y + h \cdot \text{func}(x, y)$ \\ 
        $x \leftarrow x + h$
      };
      \node [preparation, on chain, join] (inc_i) {$i \leftarrow i + 1$};
      
      \node [data, right=2cm of check_i, node distance=4.5cm] (ret) {Return $y_b$};
      \node [terminator, below=of ret] (end) {End};
    \end{scope}

    \draw [thick, -Stealth] (check_i) -- node[right] {Yes} (loop_body);
    \draw [thick, -Stealth] (check_i.east) -- node[near start, above] {No} (ret.west);
    
    \draw [thick, -Stealth] (inc_i.west) -- ++(-1.5,0) |- (check_i.west);
    
    \draw [thick, -Stealth] (ret.south) -- (end.north);
    \draw [thick, -Stealth] (start.east) -- (input.west);
  \end{tikzpicture}
  \captionof{scheme}{Блок-схема функции вычисления по равномерной сетке.}
\end{center}

\vspace*{\fill}

\clearpage

\vspace*{\fill}

\begin{center}
  \begin{tikzpicture}[gost flowchart]
    \begin{scope}[start chain=flow going below, node distance=8mm]
      \node [data, on chain, join] (input) {Input: $f, \varepsilon, a, y_a, b, M$};
      \node [terminator, left= of input] (start) {Start};
      
      \node [preparation, on chain, join] (init_n) {$n \leftarrow 1$};
      \node [predefined process, on chain, join] (calc_first) {$y_{prev} \leftarrow \text{calcForGrid}(f, y_a, a, b, n)$};
      
      \node [decision, on chain, join] (check_max) {$n < M$?};
      
      \node [preparation, on chain] (mult_n) {$n \leftarrow 2n$};
      \node [predefined process, on chain, join] (calc_curr) {$y_{curr} \leftarrow \text{calcForGrid}(f, y_a, a, b, n)$};
      
      \node [decision, on chain, join] (check_eps) {$|y_{prev} - y_{curr}| \le \varepsilon$?};
      
      \node [preparation, on chain] (update_prev) {$y_{prev} \leftarrow y_{curr}$};
      
      \node [data, right=of check_max, node distance=4.5cm] (ret_prev) {Return $y_{prev}$};
      \node [data, right=2cm of check_eps] (ret_curr) {Return $y_{curr}$};
      \node [terminator, above=of ret_curr, yshift=1.5cm] (end) {End};
    \end{scope}

    \draw [thick, -Stealth] (check_max) -- node[right] {Yes} (mult_n);
    \draw [thick, -Stealth] (check_max.east) -- node[near start, above] {No} (ret_prev.west);
    
    \draw [thick, -Stealth] (check_eps) -- node[right] {No} (update_prev);
    \draw [thick, -Stealth] (check_eps.east) -- node[near start, above] {Yes} (ret_curr.west);
    
    \draw [thick, -Stealth] (update_prev.west) -- ++(-2.5,0) |- (check_max.west);
    
    \draw [thick, -Stealth] (ret_prev.south) -- +(0,-1) -| (end.north);
    \draw [thick, -Stealth] (ret_curr.north) -- (end.south);
    \draw [thick, -Stealth] (start.east) -- (input.west);
  \end{tikzpicture}
  \captionof{scheme}{Блок-схема основного алгоритма решения методом Эйлера с~автоматическим выбором шага.}
\end{center}

\vspace*{\fill}